\documentclass[
10pt, % Main document font size
a4paper, % Paper type, use 'letterpaper' for US Letter paper
oneside, % One page layout (no page indentation)
%twoside, % Two page layout (page indentation for binding and different headers)
headinclude,footinclude, % Extra spacing for the header and footer
BCOR5mm, % Binding correction
]{scrartcl}
%%%%%%%%%%%%%%%%%%%%%%%%%%%%%%%%%%%%%%%%%
% Arsclassica Article
% Structure Specification File
%
% This file has been downloaded from:
% http://www.LaTeXTemplates.com
%
% Original author:
% Lorenzo Pantieri (http://www.lorenzopantieri.net) with extensive modifications by:
% Vel (vel@latextemplates.com)
%
% License:
% CC BY-NC-SA 3.0 (http://creativecommons.org/licenses/by-nc-sa/3.0/)
%
%%%%%%%%%%%%%%%%%%%%%%%%%%%%%%%%%%%%%%%%%

%----------------------------------------------------------------------------------------
%	REQUIRED PACKAGES
%----------------------------------------------------------------------------------------

\usepackage[
nochapters, % Turn off chapters since this is an article        
beramono, % Use the Bera Mono font for monospaced text (\texttt)
eulermath,% Use the Euler font for mathematics
pdfspacing, % Makes use of pdftex’ letter spacing capabilities via the microtype package
dottedtoc % Dotted lines leading to the page numbers in the table of contents
]{classicthesis} % The layout is based on the Classic Thesis style

\usepackage{arsclassica} % Modifies the Classic Thesis package
\usepackage{ulem}
\usepackage{graphicx}
\usepackage{float}
\usepackage{algorithm2e}
\usepackage[export]{adjustbox}
\usepackage[utf8]{inputenc}  
\usepackage[english,frenchb]{babel} 
\usepackage[T1]{fontenc} % Required for accented characters
\usepackage{multicol}
\usepackage{graphicx} % Required for including images
\graphicspath{{Figures/}} % Set the default folder for images

\usepackage{enumitem} % Required for manipulating the whitespace between and within lists

\usepackage{lipsum} % Used for inserting dummy 'Lorem ipsum' text into the template


\expandafter\def\csname ver@subfig.sty\endcsname{}


\usepackage{svg}
\usepackage{hyperref}

\usepackage{amsmath,amssymb,amsthm} % For including math equations, theorems, symbols, etc

\usepackage{varioref} % More descriptive referencing

%----------------------------------------------------------------------------------------
%	THEOREM STYLES
%---------------------------------------------------------------------------------------

\theoremstyle{definition} % Define theorem styles here based on the definition style (used for definitions and examples)
\newtheorem{definition}{Definition}

\newtheorem*{remark}{Remark}
\theoremstyle{plain} % Define theorem styles here based on the plain style (used for theorems, lemmas, propositions)
\newtheorem{theorem}{Theorème}
\newtheorem{lemma}[theorem]{Lemme}
\newtheorem{proposition}[theorem]{Proposition}

\theoremstyle{remark} % Define theorem styles here based on the remark style (used for remarks and notes)

%----------------------------------------------------------------------------------------
%	HYPERLINKS
%---------------------------------------------------------------------------------------

\hypersetup{
%draft, % Uncomment to remove all links (useful for printing in black and white)
colorlinks=true, breaklinks=true, bookmarks=true,bookmarksnumbered,
urlcolor=webbrown, linkcolor=RoyalBlue, citecolor=webgreen, % Link colors
pdftitle={}, % PDF title
pdfauthor={\textcopyright}, % PDF Author
pdfsubject={}, % PDF Subject
pdfkeywords={}, % PDF Keywords
pdfcreator={pdfLaTeX}, % PDF Creator
pdfproducer={LaTeX with hyperref and ClassicThesis} % PDF producer
}

\hyphenation{Fortran hy-phen-ation} 

\title{\normalfont{FTML Exercices 2}\\
} 
\author{\spacedlowsmallcaps{}
} 
\date{\today}
\begin{document}
\renewcommand{\sectionmark}[1]{\markright{\spacedlowsmallcaps{#1}}}
\lehead{\mbox{\llap{\small\thepage\kern1em\color{halfgray} \vline}\color{halfgray}\hspace{0.5em}\rightmark\hfil}} 
\pagestyle{scrheadings}
\maketitle 
\setcounter{tocdepth}{3} 

\tableofcontents

\section{\large\color{Blue}Ordinary least squares}

\subsubsection{Enoncé}
\label{par:enoncé}

On veut étudier la fonction objectif du problème OLS présenté lors du tp 2 afin de prouver par la suite la valeur de l'estimateur OLS. Les questions 1 et 2 peuvent être traitées indépendamment. Les définitions de la convexité et du gradient sont disponibles dans lecture\textunderscore notes.pdf (ou d'autres références de votre choix).
\\

Soit $n$ et $d\in \mathbb{N}^*$.
\\

1) Soit $X\in \mathbb{R}^{n, d}$, et $y\in \mathbb{R}^n$. Calculer le gradient de 

$$
g = \left\{
    \begin{array}{ll}
        \mathbb{R}^d \rightarrow \mathbb{R}  \\
	\theta\mapsto ||X\theta-y||^2
    \end{array}
\right.
$$

L'estimateur OLS est la valeur $\hat{\theta}$ qui minimise $g$ sur $\mathbb{R}^d$.
\\

2) On veut montrer que la fonction $g$ est convexe. Il
y a de nombreuses méthodes pour cela mais utiliser ici les étapes suivantes:

\begin{itemize}
    \item[a)] montrer que si $s : \mathbb{R}^d \rightarrow \mathbb{R}^n$ est linéaire
        et $f: \mathbb{R}^n \rightarrow \mathbb{R}$ est convexe, alors $ f\circ
        s: \mathbb{R}^d \rightarrow \mathbb{R}$ est convexe.
    \item[b)] montrer que toute norme sur $\mathbb{R}^n$ est convexe.
    \item[c)] montrer que si $w: \mathbb{R}\rightarrow \mathbb{R}$ est convexe
        croissante et $a: \mathbb{R}^n \rightarrow \mathbb{R}$ est convexe, alors
        $u = w\circ a: \mathbb{R}^n \rightarrow \mathbb{R}$ est convexe.
    \item[d)] montrer que si $u: \mathbb{R}^n \rightarrow \mathbb{R}$ est convexe, alors si $\beta\in \mathbb{R}^n$, l'application $f:x\mapsto u(x+\beta)$ est convexe.
    \item[e)] Appliquer les résultats précédents pour montrer la convexité de $g$.
\end{itemize}


3) En déduire que si $X^TX$ est inversible, l'estimateur OLS est égal à 

\begin{equation}
	\hat{ \theta}= (X^TX)^{-1}X^Ty
    \label{eq:ols}
\end{equation}

4) Facultatif: montrer que si $d<n$ et $X$ est injective, alors $X^TX$ est inversible.

\end{document}
